\chapter{Коврик и клетки}\label{ch:kovrik}

\textbf{Что вы узнете:} как устроен коврик, как узнать номер клетки и что такое четверть. В этом уроке нет программ, зато он нужен для всех следующих.

\section{Коврик 12 на 12}

Дроны летают над ковриком с метками. Дрон смотрит вниз, видит метки под собой и по ним понимает, где он находится. Коврик поделён на квадратные клетки: \textbf{12 по горизонтали и 12 по вертикали}. Сторона клетки \textbf{40\cm}, значит весь коврик 4,8\,м на 4,8\,м, всего 144 клетки с номерами от 0 до 143.

\begin{figure}[tbp]
  \centering
  % Коврик 12 x 12: номера клеток, оси, пример клетки 97
\begin{tikzpicture}[x=1.08cm, y=1.08cm, font=\sffamily]
  \foreach \r in {0,...,11}{
    \foreach \c in {0,...,11}{
      \pgfmathtruncatemacro{\n}{\c + 12*\r}
      \ifnum\c<6
        \ifnum\r<6
          \fill[accentsoft] (\c,\r) rectangle (\c+1,\r+1);
        \fi
      \fi
      \node[font=\tiny, text=ink] at (\c+0.5,\r+0.5) {\n};
    }
  }
  \fill[accent!35] (1,8) rectangle (2,9);
  \node[font=\tiny\bfseries, text=ink] at (1.5,8.5) {97};
  \draw[rule, thin] (0,0) grid (12,12);
  \draw[accent, very thick] (0,0) rectangle (6,6);
  \draw[ink, thick] (0,0) rectangle (12,12);
  \foreach \c in {0,...,11}{ \node[font=\scriptsize, text=inkgray] at (\c+0.5,-0.3) {\c}; }
  \foreach \r in {0,...,11}{ \node[font=\scriptsize, text=inkgray] at (-0.3,\r+0.5) {\r}; }
  \node[text=inkgray, font=\small] at (6,-0.85) {столбец (ось $x$)};
  \node[text=inkgray, font=\small, rotate=90] at (-0.85,6) {строка (ось $y$)};
  \node[anchor=west, text=accent, font=\small\bfseries, fill=white, inner sep=1pt] at (6.15,5.6) {четверть 6$\times$6};
  \draw[<->, inkgray] (0,12.35) -- (12,12.35) node[midway, above, font=\small] {12 клеток по 40\,см = 4,8\,м};
  \node[anchor=west, font=\scriptsize, text=ink, fill=white, inner sep=1pt] at (2.1,8.5) {97 = 1 + 12$\times$8};
\end{tikzpicture}

  \caption{Коврик 12$\times$12. Клетка 0 внизу слева, номера растут вправо и вверх. Рамкой выделена четверть 6$\times$6.}\label{fig:mat}
\end{figure}

\section{Номер клетки}

У каждой клетки есть \emph{столбец} (слева направо, от 0 до 11) и \emph{строка} (снизу вверх, от 0 до 11). Номер считается так:
\[
  \text{номер} = \text{столбец} + 12 \times \text{строка}.
\]
Клетка 97: столбец 1, строка 8, потому что $1 + 12\times 8 = 97$. Обратно: разделите номер на 12, целая часть — строка, остаток — столбец. $97 = 12\times 8 + 1$.

Клетка 0 внизу слева. Ось~$x$ идёт вправо, ось~$y$ вверх. Дрон летит в центр клетки: центр клетки 97 находится в точке $x = 40$\cm, $y = 320$\cm.

\section{Четверть}

В классе такой большой коврик обычно не разложишь. Коврик складывают в четыре, и получается квадрат 6$\times$6 клеток, то есть 2,4\,м на 2,4\,м. Это \textbf{четверть}. Большинство примеров в книге рассчитано на одну четверть.

Номера клеток всегда считаются по всему коврику. В левой нижней четверти это клетки со столбцами 0–5 и строками 0–5 (рисунок~\ref{fig:quarter}). Вот их номера по строкам: 0–5, 12–17, 24–29, 36–41, 48–53, 60–65.

\begin{figure}[tbp]
  \centering
  % Четверть 6 x 6 с маршрутом 13 -> 16 -> 52 -> 49 -> 13
\begin{tikzpicture}[x=1.45cm, y=1.45cm, font=\sffamily]
  \foreach \r in {0,...,5}{
    \foreach \c in {0,...,5}{
      \pgfmathtruncatemacro{\n}{\c + 12*\r}
      \node[font=\scriptsize, text=inkgray, anchor=north west] at (\c+0.03,\r+0.97) {\n};
    }
  }
  \draw[rule] (0,0) grid (6,6);
  \draw[ink, thick] (0,0) rectangle (6,6);
  \foreach \c in {0,...,5}{ \node[font=\small, text=inkgray] at (\c+0.5,-0.3) {\c}; }
  \foreach \r in {0,...,5}{ \node[font=\small, text=inkgray] at (-0.3,\r+0.5) {\r}; }
  \coordinate (a) at (1.5,1.5); \coordinate (b) at (4.5,1.5);
  \coordinate (c) at (4.5,4.5); \coordinate (e) at (1.5,4.5);
  \draw[accent, very thick, -{Stealth[length=3mm]}] (a) -- (b);
  \draw[accent, very thick, -{Stealth[length=3mm]}] (b) -- (c);
  \draw[accent, very thick, -{Stealth[length=3mm]}] (c) -- (e);
  \draw[accent, very thick, -{Stealth[length=3mm]}] (e) -- (a);
  \foreach \p/\t in {a/13, b/16, c/52, e/49}{
    \fill[white, draw=accent, very thick] (\p) circle (0.2);
    \node[font=\scriptsize\bfseries, text=accent] at (\p) {\t};
  }
  \node[font=\small, text=inkgray] at (3,-0.85) {40\,см в каждой клетке; вся четверть $2{,}4\times2{,}4$\,м};
\end{tikzpicture}

  \caption{Левая нижняя четверть и маршрут 13, 16, 52, 49, 13: квадрат со стороной в три клетки.}\label{fig:quarter}
\end{figure}

\section{Метки на коврике}

Метки на коврике называются AprilTag: это чёрно-белые квадраты, похожие на упрощённые QR-коды. Дрон узнаёт метку камерой и определяет по ней своё место. Позиция у дрона появляется не сразу после включения, а в полёте, примерно на высоте 30\cm. Точность около $\pm 10$–$15$\cm: поэтому в шоу не стоит планировать проходы с зазором меньше этого.

\begin{probuy}
  \item Какой номер у клетки в столбце 3 и строке 2?
  \item В каком столбце и в какой строке клетка 40?
  \item Чему равны $x$ и $y$ центра клетки 52?
  \item Сколько сантиметров между центрами клеток 13 и 16? А между 13 и 52?
  \item Сколько клеток в четверти? Какой номер у её правой верхней клетки?
  \item Какой номер у клетки прямо над клеткой 97?
\end{probuy}
\otvety{5}
